\section{Appendix}

\subsection{Prompt Formats}

All variants use the same fact and rule context and differ only in the target format. The prompt prefix is:
\begin{verbatim}
You are a careful logical reasoner.
Use only the numbered facts and rules below.
\end{verbatim}

The main task formats are:
\begin{itemize}
\item \textbf{answer-only}: \texttt{LABEL: <True|False|Unknown>}
\item \textbf{proof-only}: \texttt{MODE}, \texttt{LABEL}, \texttt{CHAIN}, and \texttt{WITNESS}, with \texttt{CHAIN} allowed to be \texttt{NONE}
\item \model{} and its ablations: the same four fields, with \texttt{CHAIN} used for either a derivation path or a \texttt{FAIL[...]} abstention tag
\end{itemize}

This shared prompt structure matters for the comparison. The main difference between variants is the supervision target, not the surrounding instruction.

\subsection{Training Hyperparameters}

All runs use QLoRA with:
\begin{itemize}
\item 4-bit NF4 loading,
\item LoRA rank 16,
\item LoRA alpha 32,
\item dropout 0.05,
\item max sequence length 512,
\item learning rate \(2\times 10^{-4}\),
\item one training epoch.
\end{itemize}

The model-specific batch settings are:
\begin{itemize}
\item \textbf{0.5B pilot}: per-device train batch 8, gradient accumulation 4
\item \textbf{7B main}: per-device train batch 2, gradient accumulation 8
\end{itemize}

The initial pilot figures use the same 4,096-example train subset, 1,000-example ID subset, and 500-example OOD subset, with seed 0. The upgraded main-track sweeps retain the same training-budget controls, but evaluate multi-seed runs on a fixed sampled test subset by holding the evaluation data seed constant. Experiments were run on NVIDIA A100 80GB GPUs.

\subsection{Verifier Details}

The verifier replays generated outputs against the benchmark's symbolic structure:
\begin{itemize}
\item For \textbf{PROVE} outputs, the predicted \texttt{CHAIN} must derive the queried literal.
\item For \textbf{REFUTE} outputs, the predicted \texttt{CHAIN} must derive the opposite literal.
\item For \textbf{ABSTAIN} outputs, the \texttt{WITNESS} must describe missing support that is absent from the theory closure, or the chain must explicitly use a \texttt{FAIL[no\_rule]} tag.
\end{itemize}

This verifier is deliberately lightweight. It does not search for an alternative proof on the model's behalf. It only checks the consistency of the model's own evidence claim.

\subsection{Experiment Automation}

Main-track sweeps are launched with a single manifest-driven script:
\begin{verbatim}
python scripts/run_main_track_suite.py
\end{verbatim}
This launcher keeps per-run metadata, GPU-specific logs, and deterministic output names. After new result JSONs are produced, summaries, plots, LaTeX tables, and unknown-behavior reports are refreshed with:
\begin{verbatim}
python scripts/refresh_artifacts.py --compile-paper
\end{verbatim}

\subsection{Claim-to-Evidence Map}

Table~\ref{tab:claims-map} links each headline claim in the paper to one exact result artifact and states the scope of that claim. We use this table as a guardrail against vague summary language.

\begin{table}[t]
\centering
\small
\begin{tabular}{p{0.36\linewidth}p{0.34\linewidth}p{0.22\linewidth}}
\toprule
Claim & Supporting Evidence & Scope \\
\midrule
At 7B and 4k train, ProCo greatly improves in-domain joint accuracy over proof-only. & 75.0 $\pm$ 0.7 vs 39.9 $\pm$ 0.6 (Table~\ref{tab:seeded-results}) & Qwen2.5-7B, depth-3ext-NatLang/test, subset 4000, 3 seeds. \\
At 7B and 4k train, the same gain remains on depth-OOD. & 55.0 $\pm$ 0.7 vs 26.5 $\pm$ 1.1 (Table~\ref{tab:seeded-results}) & Qwen2.5-7B, depth-5/test, subset 4000, 3 seeds. \\
The unknown-case gain comes from explanation quality rather than abstention rate. & ID predicted unknown: 1525.3 vs 1544.0; ID faithful unknown: 1444.3 vs 0.0. Depth-OOD faithful unknown: 1169.0 vs 0.0 (Table~\ref{tab:unknown-behavior-maintrack}) & Gold-unknown slices on the fixed 7B test subset. \\
With full training, the in-domain raw-accuracy gap is nearly gone. & ProCo acc. 98.2 vs answer-only 98.5; ProCo joint 96.8 (Table~\ref{tab:scaling-results}) & Qwen2.5-7B, depth-3ext-NatLang/test, subset 4000, train=112,062. \\
With full training, ProCo slightly surpasses answer-only on depth-5 subset accuracy. & ProCo acc. 93.8 vs answer-only 93.1 (Table~\ref{tab:scaling-results}) & Qwen2.5-7B, depth-5/test, subset 4000, train=112,062. \\
\bottomrule
\end{tabular}
\caption{Claim-to-evidence map for the main paper. Each row states one headline claim, the exact number or table that supports it, and the scope of that evidence.}
\label{tab:claims-map}
\end{table}

\subsection{Per-Class Breakdown on the 7B Seeded Suite}

Table~\ref{tab:per-class-seeded} reports per-class F1 on the fixed 4k seeded 7B suite. At this training budget, answer-only remains the strongest raw classifier by classwise F1. The main advantage of \model{} at this scale appears instead in faithful evidence and joint accuracy.

\begin{table}[t]
\centering
\small
\begin{tabular}{llccc}
\toprule
Domain & Variant & True F1 & False F1 & Unknown F1 \\
\midrule
ID & answer-only & 91.2 $\pm$ 0.1 & 88.5 $\pm$ 0.3 & 87.7 $\pm$ 0.4 \\
ID & proof-only & 90.3 $\pm$ 0.3 & 87.2 $\pm$ 0.3 & 86.7 $\pm$ 0.3 \\
ID & ProCo & 89.9 $\pm$ 0.3 & 87.0 $\pm$ 0.2 & 86.0 $\pm$ 0.6 \\
\midrule
Depth-OOD & answer-only & 82.8 $\pm$ 1.6 & 82.4 $\pm$ 1.0 & 76.3 $\pm$ 0.6 \\
Depth-OOD & proof-only & 80.2 $\pm$ 2.2 & 76.7 $\pm$ 0.6 & 73.0 $\pm$ 0.5 \\
Depth-OOD & ProCo & 80.5 $\pm$ 2.2 & 76.9 $\pm$ 0.8 & 73.3 $\pm$ 1.0 \\
\bottomrule
\end{tabular}
\caption{Per-class F1 on the fixed 7B seeded suite. This table helps separate raw class prediction from evidence faithfulness.}
\label{tab:per-class-seeded}
\end{table}

\subsection{Metric Interpretation}

Joint accuracy is defined as correct label plus faithful evidence. This gives answer-only joint accuracy 0 by construction, because answer-only systems do not emit evidence to verify. We therefore read joint accuracy together with raw accuracy throughout the paper: answer-only measures label prediction alone, while proof-only and \model{} are asked to produce checkable reasoning.

\subsection{Claim Audit Artifact}

Before any submission draft leaves the repository, we run:
\begin{verbatim}
python scripts/check_paper_claims.py
\end{verbatim}
This script checks the headline numbers in the abstract, introduction, and discussion against the aggregated result files and writes the audit to \texttt{docs/PAPER\_CLAIM\_CHECK\_2026-04-26.md}. The current audit passes.
